\section{Estado del arte}
\label{sec:arte}
Los apartados siguientes establecen el recorrido de la revisión. Se presentan los métodos estudiados y se señalan los contenidos pendientes de desarrollo.

\subsection{Evaluación de calidad de imagen (IQA)}
\label{sec:iqa}
\subsubsection{Conceptos y factores de degradación}
\pendiente{Definir Image Quality Assessment (IQA) y describir iluminación, desenfoque, ruido, resolución y compresión, con referencias.}
\subsubsection{Métodos con y sin referencia}
\pendiente{Distinguir referencia completa (FR-IQA), referencia reducida (RR-IQA) y ausencia de referencia (NR-IQA); aclarar qué información necesita cada familia.}
\subsubsection{De la calidad perceptual a la calidad biométrica facial}
\pendiente{Introducir pose, expresión, oclusiones y condiciones de captura. Explicar cómo se verifica la relación entre calidad estimada y rendimiento biométrico.}

\subsection{Modelos de visión-lenguaje (VLMs)}
\subsubsection{Representaciones visuales y textuales}
\pendiente{Presentar el aprendizaje multimodal y la alineación entre imagen y texto.}
\subsubsection{Arquitecturas contrastivas y generativas}
\pendiente{Revisar CLIP, Qwen-VL y LLaVA a partir de sus artículos originales, indicando versión y diferencias relevantes para FIQA.}
\subsubsection{Estrategias de adaptación}
\pendiente{Definir zero-shot, few-shot y fine-tuning. Registrar ejemplos, entrenamiento y datos utilizados en cada trabajo.}

\subsection{Aprendizaje profundo aplicado a IQA y FIQA}
\subsubsection{De métricas tradicionales a redes profundas}
\pendiente{Resumir la evolución hacia redes convolucionales (CNN) y transformers sin desviar el foco de la utilidad biométrica.}
\subsubsection{Métodos especializados de FIQA}
Las representaciones mediante \textit{embeddings} permiten estudiar la calidad de las imágenes faciales en relación con su utilidad para el reconocimiento biométrico. Un ejemplo es SER-FIQ (Terhorst et al., 2020)~\cite{Terhorst_2020_CVPR}, que propone una estimación de calidad sin etiquetas de calidad, basada en la robustez de las representaciones faciales. Para ello, cuantifica la dispersión de múltiples \textit{embeddings} obtenidos de una misma imagen mediante Monte Carlo Dropout. En su evaluación experimental, realizada con los modelos de reconocimiento facial FaceNet y ArcFace sobre los conjuntos de datos Adience y LFW, el método supera a distintas aproximaciones académicas y comerciales en la mayoría de los escenarios analizados. En particular, muestra una mayor capacidad para identificar y descartar imágenes de baja calidad, reduciendo los errores de verificación facial conforme aumenta la proporción de imágenes rechazadas. Estos resultados sugieren que la estabilidad de los \textit{embeddings} aporta información relevante para estimar la utilidad biométrica de una imagen y constituye una base para explorar otros métodos de evaluación de calidad basados en representaciones profundas.

SDD-FIQA (Ou et al., 2021)~\cite{Ou_2021_CVPR} lleva un paso más allá el uso de \textit{embeddings} y el enfoque de evaluación de calidad no supervisada, incorporando tanto la similitud entre imágenes de una misma identidad (intraclase) como la similitud entre imágenes de identidades diferentes (interclase). Para ello, propone medir la separación entre ambas distribuciones mediante la distancia de Wasserstein, considerando que cuanto mayor sea esta separación, mayor será la calidad biométrica de la imagen. A partir de estas distancias, genera pseudoetiquetas de calidad que permiten entrenar una red de regresión encargada de predecir la calidad de nuevas imágenes faciales.

\pendiente{Revisar FaceQAN, ViT-FIQA y eDifFIQA. Para cada método, documentar señal de calidad, entrenamiento y dependencia del reconocedor facial.}
\subsubsection{Alcance de las puntuaciones escalares}
\pendiente{Analizar qué permite evaluar una puntuación y qué información adicional se necesita para explicar una degradación.}

\subsection{VLMs aplicados a IQA y FIQA}
\subsubsection{Estimación de calidad y descripción de degradaciones}
\pendiente{Revisar Q-Align y CLIB-FIQA, verificando sus objetivos, arquitecturas y protocolos. Separar evidencia de IQA general y de FIQA.}
\subsubsection{VLMs generalistas en escenarios zero-shot}
El contexto propone como punto de partida el artículo \textit{Employing Vision-Language Models for Face Image Quality Assessment}, atribuido a Sarıtaş, Onaran, Štruc y Ekenel (2026). Sus metadatos y su contenido quedan pendientes de verificación antes de incorporarlo como referencia científica o describir sus resultados.
\pendiente{Localizar y leer el artículo indicado en el contexto:\newline
\url{https://arxiv.org/abs/2605.17489}.\newline
Contrastar el código asociado:\newline
\url{https://github.com/ThEnded32/VLM4FIQA}.}
\subsubsection{Tipos de salida y sistemas híbridos}
\pendiente{Comparar puntuaciones, atributos, explicaciones y combinaciones de FIQA especializado con VLMs. Distinguir consistencia textual y exactitud de las explicaciones.}
